\chapter{2021 年试题（当年科目代码 818）}

\begin{tishi}
\textbf{关于科目代码：}2021 年西华大学该科目的\textbf{考试科目代码为 818}（名称为"工程流体力学"），
与现行 818 一致（已用 \texttt{pdftoppm} 渲染原卷第 1、3、4、5 页目视复核页脚，
均为"考试科目代码及名称：818 工程流体力学"）。本章\textbf{不涉及旧大纲（X1 明渠流与堰流、
X2 气体动力学、X3 流体机械）取分的题目}，全部内容落在赵琴教材第 1--8 章范围内。

\textbf{试卷说明（原卷原文）：}试题附在答卷内交回，答案必须写在答题纸上，写在试题纸上无效。

\textbf{结构：}一、单项选择题 10 小题 \(\times\) 3 分 \(=\) 30 分；
二、判断题 10 小题 \(\times\) 2 分 \(=\) 20 分；
三、填空题\textbf{共 10 空，每空 3 分} \(=\) 30 分；
四、简答题 3 小题 \(\times\) 4 分 \(=\) 12 分；
五、计算题 6 小题 \(=\) 58 分。满分 150 分。

\textbf{注意：}原卷正文因答题纸装订（"试题附在答卷内交回"）而\textbf{题号发生串接}：
试卷上"四．简答题"之前的填空部分被排为两处\textbf{未标序号的"三、填空题"}区段，
且"四．简答题"的\textbf{节次上出现了重号}。此处\textbf{按内容与分值重排}为
"三、填空题"（10 空）、"四、简答题"（12 分）、"五、计算题"（58 分），
\textbf{题干文字一律照原卷转写}，请对照原卷影印核实。

\textbf{图示：}本题全卷 8 幅插图均未用代码绘制，凡有图之题在题干后以
\tuyi{} 给出文字图示说明，并注明"原卷影印见本册附录"。
\end{tishi}

\section{一、单项选择题（本题共 10 小题，每小题 3 分，小计 30 分）}

\begin{ti}
  \item 压力体内（\quad）。\hfill\xstarfull{5}\yiju{E5 题库计算 1（弧形闸门 $F_x,F_z$ 与力矩）；E1 slide33}\nandu{易}
  \begin{choices}
    \item 必定充满液体
    \item 肯定不会有液体
    \item 至少部分有液体
    \item 可能有液体，也可能无液体
  \end{choices}

  \item 重力场中，静止液体所受的单位质量力为（\quad）。\hfill\xstarfull{5}\yiju{E4 2022 单选 3（测压管水头线坡度）、单选 7（虹吸管）；E1 slide33"计算题基础入门"}\nandu{易}
  \begin{choices}
    \item $f_x=f_y=f_z=0$
    \item $f_x=0,\ f_y=-g,\ f_z=0$
    \item $f_x=-g,\ f_y=f_z=0$
    \item $f_x=f_y=0,\ f_z=-g$
  \end{choices}

  \item 动量方程中，\(F\) 表示作用在控制体内流体上的力是（\quad）。\hfill\xstarfull{5}\yiju{E4 2015 计算（水平弯管求 $F_x,F_y$）；E4 2026 计算 3（教材习题 4.12）；E1 slide16}\nandu{易}
  \begin{choices}
    \item 总质量力
    \item 总表面力
    \item 合外力
    \item 总压力
  \end{choices}

  \item 已知不可压缩流体平面流动的流速场为 \(u_x=xt\)，\(u_y=-yt\)，
  那么在时刻 \(t=1\,\mathrm{s}\) 时点 \(A(1,\,2)\) 处流体质点的加速度为（\quad）。
  \hfill\xstarfull{5}\yiju{E4 各年选择/填空（质点导数与加速度）；E1 slide33"公式较多，重点记忆"}\nandu{中}
  \begin{choices}
    \item $a_x=0,\ a_y=2$
    \item $a_x=2,\ a_y=0$
    \item $a_x=2,\ a_y=4$
    \item $a_x=4,\ a_y=2$
  \end{choices}

  \item 已知不可压缩流体的流速场为 \(u_x=f(y)\)，\(u_y=f(x)\)，\(u_z=0\)，
  则该流动为（\quad）。\hfill\xstarfull{4}\yiju{E4 2022 单选 4（方管进压缩机求出口密度与质量流量）}\nandu{易}
  \begin{choices}
    \item 一元流
    \item 二元流
    \item 三元流
    \item 均匀流
  \end{choices}

  \item 应用动量方程求流体对固体壁面的作用力时，进、出口截面上的压强应使用（\quad）。
  \hfill\xstarfull{5}\yiju{E4 2015 计算（水平弯管求 $F_x,F_y$）；E4 2026 计算 3（教材习题 4.12）；E1 slide16}\nandu{中}
  \begin{choices}
    \item 绝对压强
    \item 相对压强
    \item 大气压
    \item 真空度
  \end{choices}

  \item 管道中液体的雷诺数与（\quad）无关。\hfill\xstarfull{5}\yiju{E4 2022 单选 9（两圆管流态判别）；E4 2026 简答 3（层流与雷诺数关系）}\nandu{易}
  \begin{choices}
    \item 温度
    \item 管径
    \item 管长
    \item 流速
  \end{choices}

  \item 阻力平方区的沿程水头损失与管截面上的平均流速的（\quad）成正比。\hfill\xstarfull{3}\yiju{E6 教材 5.4；E2 slide5}\nandu{中}
  \begin{choices}
    \item \(\dfrac{1}{2}\) 次方
    \item \(1.75\) 次方
    \item 二次方
    \item 一次方
  \end{choices}

  \item 用于判别层流和紊流的是（\quad）。\hfill\xstarfull{5}\yiju{E4 2022 单选 9（两圆管流态判别）；E4 2026 简答 3（层流与雷诺数关系）}\nandu{易}
  \begin{choices}
    \item 上临界雷诺数
    \item 下临界雷诺数
    \item 上、下临界雷诺数的代数平均数
    \item 上、下临界雷诺数的几何平均数
  \end{choices}

  \item 两水箱之间用三根不同直径相同长度的水平管道 1，2，3 相连接，如图 1，
  三管沿程阻力系数相等，则下面关系式正确的是（\quad）。\hfill\xstarfull{4}\yiju{E1 slide33；E6 教材 5.7}\nandu{难}
  \begin{choices}
    \item $\dfrac{Q_1^2}{d_1^4}=\dfrac{Q_2^2}{d_2^4}=\dfrac{Q_3^2}{d_3^4}$
    \item $\dfrac{Q_1^2}{d_1^5}=\dfrac{Q_2^2}{d_2^5}=\dfrac{Q_3^2}{d_3^5}$
    \item $\dfrac{Q_1}{d_1^2}=\dfrac{Q_2}{d_2^2}=\dfrac{Q_3}{d_3^2}$
    \item $\dfrac{Q_1^2}{d_1^3}=\dfrac{Q_2^2}{d_2^3}=\dfrac{Q_3^2}{d_3^3}$
  \end{choices}
  \tuyi{原卷图 1：两个敞口（或封闭）水箱之间用\textbf{三根互相平行、长度相同、
  直径依次不同}的水平圆管道 1、2、3 连通，三管的直径分别为 \(d_1>d_2>d_3\)，
  各管的流量分别为 \(Q_1\)、\(Q_2\)、\(Q_3\)；三管为并联关系，
  故三管的沿程水头损失相等。（原卷影印见本册附录）}
\end{ti}

\begin{tishi}
\textbf{第 2 题：}原卷四个选项为竖排的公式，OCR 只拾得零散符号
（作"\(f_x=f_y=f_z=0\)""\(f_x=0,f_y=-g,f_z=0\)""\(f_x=-g,f_y=f_z=0\)""\(f_x=f_y=0,f_z=-g\)"），
其中\textbf{哪个分量等于零、哪个分量等于 \(-g\)} 的字形已被 OCR 合并；
此处照 OCR 拾得的四个式子\textbf{字面转写，未改动}，
\textbf{选项与公式的对应顺序请对照原卷影印核实}。

\textbf{第 4 题：}原卷选项为竖排公式，OCR 正文中混入了页码外的演算字迹
（作"\(=x+xt\)""\(=-y+(-yt)(-t)\)"），那是原卷\textbf{手写笔算的痕迹，不是选项文字}，
此处已剔除；四个选项照原卷转写为 \(a_x,a_y\) 的四组取值。

\textbf{第 10 题：}原卷四个选项为竖排分式，OCR 丢失了分子分母的对应关系；
据原卷影印，分子均为 \(Q_i\) 的平方、分母均为 \(d_i\) 的若干次幂，
四个选项的指数依次为 \(4\)、\(5\)、\(2\)、\(3\)。
此处照影印转写，若与您手中的原卷不同，请以原卷为准。
\end{tishi}

\section{二、判断题（正确的填"$\surd$"，错误的填"$\times$"，本题共 10 小题，每小题 2 分，小计 20 分）}

\begin{ti}
  \item 流体静止时不考虑流体内的切应力。\hfill\xstarfull{5}\yiju{E4 2026 计算 1（牛顿流体斜板静压强）；E5 题库选择 2（汽油＝牛顿流体）}\nandu{易}

  \item 均匀流的迁移（位变）加速度等于零。\hfill\xstarfull{3}\yiju{E1 slide16 简答（均匀流的特点）}\nandu{易}

  \item 凡是满管流流动，任何断面上的压强均大于大气的压强。\hfill\xstarfull{4}\yiju{E4 2018 选择 3（U 形水银测压计）；E4 2015 填空（U 形压差计）}\nandu{中}

  \item 有压管道中，流体的机械能损失可以根据流动边界形状的变化情况分成沿程损失和局部损失两类。
  \hfill\xstarfull{5}\yiju{E1 slide16 简答"水头损失的分类及定义"；E1 slide33；E2 slide5}\nandu{易}

  \item 圆管层流的沿程阻力系数仅与雷诺数有关。\hfill\xstarfull{5}\yiju{E4 2022 单选 3；E5 题库选择 10（圆管与方管沿程损失比 0.886）}\nandu{易}

  \item 当管道出口与一充分大的容器相连时，液体在出口处的局部阻力系数等于零。
  \hfill\xstarfull{4}\yiju{E4 2015 计算 3（突然缩小水头损失）；E1 slide16}\nandu{中}

  \item 按短管计算时管道的速度水头不能略去。\hfill\xstarfull{4}\yiju{E4 2022 单选 7（虹吸管）；E5 题库选择 8、14}\nandu{易}

  \item 水箱通过孔口泄水，当水位下降时的非恒定流的泄水时间，小于水位不变时恒定泄放同体积水体所需的时间。
  \hfill\xstarfull{5}\yiju{E1 slide33；E2 slide5；E4 多卷计算}\nandu{中}

  \item 所谓水力粗糙管是指内壁面粗糙度很大的管道。\hfill\xstarfull{3}\yiju{E6 教材 5.4；E2 slide5}\nandu{中}

  \item 如图 2，流量为 \(Q\)、平均流速为 \(v\) 的射流，冲击直立平板后分成两股，
  一股沿板面直泻而下，流量为 \(Q_1\)，另一股以倾角 \(\alpha\) 射出，流量为 \(Q_2\)，
  两股分流的平均速度均等于 \(v\)。若不计摩擦力和重力影响，
  倾角 \(\alpha\) 的计算式为 \(\cos\alpha=\dfrac{Q_2}{Q}\)。\hfill\xstarfull{5}\yiju{E4 2015 计算（水平弯管求 $F_x,F_y$）；E4 2026 计算 3（教材习题 4.12）；E1 slide16}\nandu{难}
  \tuyi{原卷图 2：一股水平射流（流量 \(Q\)、平均流速 \(v\)）沿水平方向冲击一块竖直平板；
  取平板为界，射流分成两股：一股沿板面\textbf{竖直向下}流走（流量 \(Q_1\)、流速 \(v\)），
  另一股沿与铅垂方向成倾角 \(\alpha\) 的方向斜向射出（流量 \(Q_2\)、流速 \(v\)）；
  断面 0--0 取在射流进口，断面 1--1 取在下股出口，断面 2--2 取在斜射出口。
  （原卷影印见本册附录）}
\end{ti}

\begin{tishi}
\textbf{第 10 题：}原卷该式排版分两行且分子处的字形 OCR 拾得略模糊
（曾作"\(\cos\alpha=\dfrac{Q_2}{Q_2}\)"，分母明显为 \(Q\)）；
此处按流体力学常识与影印件写成 \(\cos\alpha=\dfrac{Q_2}{Q}\)。
若原卷确为其他写法，请以原卷为准。
\end{tishi}

\section{三、填空题（本题共 10 空，每空 3 分，小计 30 分）}

\begin{ti}
  \item 凡是遵守牛顿内摩擦定律的流体称为 \underline{\hspace{3cm}}。\hfill\xstarfull{4}\yiju{E5 题库选择 2（汽油＝牛顿流体）}\nandu{易}

  \item 静止流体中任一点上不同方向的静压强大小 \underline{\hspace{3cm}}。\hfill\xstarfull{4}\yiju{E2 slide4；E5 概念题高频}\nandu{易}

  \item 如图 3 所示，以 U 型管测量 \(A\) 处水压强，已知 \(h_1\)、\(h_2\)、水的重度 \(\gamma\)
  以及 U 型管中水银的重度 \(\gamma_{\mathrm{Hg}}\)，则 \(A\) 处绝对压强等于
  \underline{\hspace{3cm}}。
  \hfill\xstarfull{4}\yiju{E4 2018 选择 3（U 形水银测压计）；E4 2015 填空（U 形压差计）}\nandu{难}
  \tuyi{原卷图 3：一容器（内盛水，标注 \(A\) 点）右侧壁接出一根倒 U 型测压管；
  管内左侧自上而下先为水柱 \(h_1\)，其下为水银柱 \(h_2\)，管子右支开口通大气（标注 \(p_a\)）。
  图中 \(h_1\) 为 \(A\) 点所在水平面至水与水银分界面的铅垂距离，
  \(h_2\) 为该分界面至 U 管底（水银最低处）的铅垂距离。（原卷影印见本册附录）}

  \item 压力体可分为实压力体和虚压力体。如果压力体与液体位于受压面异侧，
  称为 \underline{\hspace{2.5cm}}，垂直分力 \underline{\hspace{2.5cm}}（填"向上"或"向下"）。
  \hfill\xstarfull{5}\yiju{E5 题库计算 1（弧形闸门 $F_x,F_z$ 与力矩）；E1 slide33}\nandu{中}

  \item \underline{\hspace{2.5cm}}（填"系统"或"控制体"）内所含流体质量保持不变。
  \hfill\xstarfull{5}\yiju{E4 2022 单选 4；E5 题库}\nandu{易}

  \item 根据流体的运动要素和空间坐标数量的关系，可将流动分为一维、二维和三维流动。
  管道和渠道流动，流动方向的尺寸远大于横向尺寸，可将流动视为 \underline{\hspace{3cm}}。
  \hfill\xstarfull{4}\yiju{E4 2022 单选 4（方管进压缩机求出口密度与质量流量）}\nandu{易}

  \item 流体运动形式中的 \underline{\hspace{3cm}} 运动是刚体不具有的。\hfill\xstarfull{4}\yiju{E5 题库选择 13（涡量为零）；E6 教材 3.4}\nandu{中}

  \item 已知平面不可压缩流体速度势函数为 \(\psi=2xy-y\)，
  则流体的速度分量 \(u_x=\) \underline{\hspace{2.5cm}}，
  \(u_y=\) \underline{\hspace{2.5cm}}。\hfill\xstarfull{5}\yiju{E4 2015 计算（流函数求速度与流线）；E4 2026 计算 4（证明势函数存在性）}\nandu{难}
\end{ti}

\begin{tishi}
\textbf{第 3 题：}原卷该空下方的版面上另有"Pa"字样与一处下划线；
据影印件，"Pa"是所给压强单位。按原卷影印，\(A\) 处绝对压强为
\(p_A=p_a+\gamma h_1-\gamma_{\mathrm{Hg}}h_2\)（\(p_a\) 为大气压强）。
\textbf{原卷在此处究竟是几空、第二个空要求填什么，OCR 无法判读
（可能只此一空，也可能是"绝对压强"与"相对压强"两空）}，
此处照影印排为\textbf{一个空}，请对照原卷影印核实。

\textbf{第 4 题：}原卷"如果压力体与液体位于受压面异侧，称为 \underline{\hspace{1.5cm}}"处的
填空线在 OCR 中\textbf{未被拾得}（OCR 只留下下一行的"垂直分力"与"（填'向上'或'向下'）"）；
据影印件该处确有下划线，故本题按\textbf{两空}（压力体名称 + 垂直分力方向）转写。
另：原卷题干为竖排两行，OCR 作"称为"/"垂直分力"/"（填...）"三行，
语序应为"……称为 \underline{\hspace{1cm}}，垂直分力 \underline{\hspace{1cm}}（填……）"。

\textbf{第 8 题：}原卷排的是速度势函数 \(\psi=2xy-y\) 并问两个速度分量；
按不可压缩流体速度势与速度的关系，\(\psi\) 实为\textbf{流函数}符号，
但\textbf{原卷字面即写 \(\psi=2xy-y\) 与该二空，此处照原卷转写、未改符号}，
请对照原卷影印核实。另 OCR 把两空都拾成"\(u=\)"（第二个应为 \(u_y\)），此处按影印补为 \(u_x\)、\(u_y\)。

\textbf{填空题空数说明：}原卷大题题头为"本题共 10 空，每空 3 分，小计 30 分"，
而按上文点数，"第 3 题"两个空、"第 4 题"两个空、"第 5--7 题"各一空、
"第 8 题"两个空，合计恰为 10 空，与题头一致；
其中"第 3 题"第二空的内容 OCR 未能读全，请对照原卷影印核实。
\end{tishi}

\section{四、简答题（本题共 3 小题，每小题 4 分，小计 12 分）}

\begin{ti}
  \item 重力场中静止流体的等压面是什么形状？为什么？\hfill\xstarfull{4}\yiju{E4 2022 单选 5（$dp$ 形式）；E5 题库选择 3（体积力有势）}\nandu{易}

  \item 尼古拉兹人工粗糙管实验曲线图中 \(\lambda\) 仅和 \(\Rey\) 有关是哪些区域，
  \(\lambda\) 仅和 \(\Delta/d\) 有关的又是哪些区域？\hfill\xstarfull{3}\yiju{E6 教材 5.4；E2 slide5}\nandu{中}

  \item 开展模型实验设计时，对于有压管流和明渠流动分别应选择哪个相似准则？
  \hfill\xstarfull{4}\yiju{E4 2014 选择 10；E4 2016 填空 4；E4 2019 判断 2；E4 2022 单选 9；E5 题库选择 12}\nandu{中}
\end{ti}

\section{五、计算题（本题共 6 小题，小计 58 分）}

\begin{ti}
  \item 如图 4，扇形闸门，中心角 \(\alpha=45^\circ\)，宽度 \(B=1\,\mathrm{m}\)（垂直于图面），
  可以绕铰链 \(C\) 旋转，用以蓄水或泄水。水深 \(H=3\,\mathrm{m}\)，
  求作用于此闸门上的静水总压力的大小和方向。（10 分）\hfill\xstarfull{5}\yiju{E5 题库计算 1（弧形闸门 $F_x,F_z$ 与力矩）；E1 slide33}\nandu{难}
  \tuyi{原卷图 4：一扇形（圆弧面）闸门，圆弧的中心 \(C\) 位于闸底水平面上，
  圆弧半径 \(r\) 由 \(C\) 指向闸门顶部 \(b\)；闸门顶端 \(b\) 与圆心 \(C\) 之间的中心角
  \(\alpha=45^\circ\)，闸门底端 \(a\) 在圆心 \(C\) 的正上方、
  即 \(aC\) 为竖直面；上游水位高程标注为 \(H\)（水深 \(H=3\,\mathrm{m}\)），
  水面位于闸门顶缘附近。水流作用在圆弧面的凹侧（上游侧），
  闸门可绕铰链 \(C\) 转动。（原卷影印见本册附录）}

  \item 如图 5，水从敞口水池沿一截面有变化的管路排出，若质量流量
  \(q_m=15\,\mathrm{kg/s}\)，\(d_1=100\,\mathrm{mm}\)，\(d_2=75\,\mathrm{mm}\)，
  不计损失，试求所需的水头 \(H\) 以及第二管段中央 \(M\) 点的相对压强。（8 分）
  \hfill\xstarfull{5}\yiju{E1 slide16"计算题：伯努利方程"；E1 slide33 与 E2 slide4"第四章＝考试计算题重点"}\nandu{中}
  \tuyi{原卷图 5：一敞口水池，池中水面高程距管道出口轴线的铅垂高度为 \(H\)；
  池底接出一段水平管路，管路先为直径 \(d_1=100\,\mathrm{mm}\) 的直管，
  经$\;$渐缩段后变为直径 \(d_2=75\,\mathrm{mm}\) 的管段，末端开放、水射入大气；
  在 \(d_2\) 管段中段标注 \(M\) 点（该点位于管轴线上）。（原卷影印见本册附录）}

  \item 如图 6，水流经水平弯管流入大气，已知 \(d_1=100\,\mathrm{mm}\)，\(d_2=75\,\mathrm{mm}\)，
  \(v_1=1.5\,\mathrm{m/s}\)，\(\theta=30^\circ\)。若不计水头损失，
  试求水流对弯管的作用力 \(F_x\)、\(F_y\)。（14 分）\hfill\xstarfull{5}\yiju{E4 2015 计算（水平弯管求 $F_x,F_y$）；E4 2026 计算 3（教材习题 4.12）；E1 slide16}\nandu{难}
  \tuyi{原卷图 6：一水平放置的渐缩弯管，进口断面 1--1 的直径 \(d_1=100\,\mathrm{mm}\)、
  进口流速 \(v_1=1.5\,\mathrm{m/s}\) 沿水平方向；出口断面 2--2 的直径 \(d_2=75\,\mathrm{mm}\)，
  出口方向相对进口方向偏转 \(\theta=30^\circ\)（管轴线在水平面内折转 \(30^\circ\)），
  出口射入大气。坐标系取 \(x\) 沿进口流速方向、\(y\) 在水平面内与 \(x\) 垂直。
  （原卷影印见本册附录）}

  \item 如图 7，钢管输水，流量为 \(52.5\,\mathrm{L/s}\)，管长为 \(300\,\mathrm{m}\)，
  管径为 \(0.2\,\mathrm{m}\)，水头为 \(20\,\mathrm{m}\)，已知沿程阻力系数为 \(0.025\)，
  局部水头损失按沿程水头损失的 \(10\%\) 计。问水塔水面要比管道出口高多少？（8 分）
  \hfill\xstarfull{5}\yiju{E5 题库计算 2（水箱垂直管道出流 $Q$ 与管长 $l$）；E1 slide33}\nandu{中}
  \tuyi{原卷图 7：一座水塔（水面标高恒定），自塔底接出一根长 \(l=300\,\mathrm{m}\)、
  管径 \(d=0.2\,\mathrm{m}\) 的钢管路，管路在中途转弯后倾斜向下、
  末端出口为水平方向自由射入大气；图中标注水塔水面至管道出口的铅垂高度为 \(H\)。
  （原卷影印见本册附录）}

  \item 有一平板，宽 \(0.5\,\mathrm{m}\)，长 \(0.8\,\mathrm{m}\)，顺流放置于水中，
  已知平板与水流的相对速度为 \(2\,\mathrm{m/s}\)，水的运动黏度为 \(10^{-6}\,\mathrm{m^2/s}\)，
  水的密度为 \(10^3\,\mathrm{kg/m^3}\)，转捩点处 \(\Rey_{x,\mathrm{cr}}=5\times10^5\)，
  求平板尾端边界层厚度以及平板两面的摩擦阻力。
  （层流边界层厚度 \(\dfrac{\delta}{x}=\dfrac{5.48}{\sqrt{\Rey_x}}\)；
  湍流边界层厚度 \(\dfrac{\delta}{x}=\dfrac{0.38}{\sqrt[5]{\Rey_x}}\)；
  层流边界层摩阻系数 \(C_f=\dfrac{1.328}{\sqrt{\Rey_l}}\)；
  湍流边界层摩阻系数 \(C_f=\dfrac{0.074}{\sqrt[5]{\Rey_l}}\)）（10 分）
  \hfill\xstarfull{2}\yiju{E4 2015 计算 9（平板摩擦阻力比 $F_1/F_2$）}\nandu{难}
  \tuyi{原卷图：一块宽 \(0.5\,\mathrm{m}\)、长 \(0.8\,\mathrm{m}\) 的薄平板顺流平放于水中，
  水流沿板长方向以 \(v=2\,\mathrm{m/s}\) 的相对速度掠过平板的上、下两面；
  板前缘为层流边界层起点，至转捩点 \(\Rey_{x,\mathrm{cr}}=5\times10^5\) 后为湍流边界层。
  （原卷影印见本册附录）}

  \item 如图 8，采用长度比尺 \(k_l=0.05\) 的模型，做弧形闸门闸下泄流实验，
  现测得模型下游收缩断面的平均流速 \(v_m=2\,\mathrm{m/s}\)，流量 \(Q_m=35\,\mathrm{L/s}\)，
  水流作用在闸门上的总压力 \(P_m=40\,\mathrm{N}\)，
  试求：原型收缩断面的平均速度 \(v_p\)、流量 \(Q_p\) 和闸门上的总压力 \(P_p\)。（8 分）
  \hfill\xstarfull{3}\yiju{E4 2014 计算 8（溢流坝模型相似）；E5 题库选择 15}\nandu{中}
  \tuyi{原卷图 8：一弧形闸门闸下泄流模型，上游为水位恒定的水库，
  弧形闸门底部开启泄流，闸下游收缩断面处水流厚度较小，
  下游渠底以下为不透水地基（图中以斜影线表示）；
  模型几何长度比尺 \(k_l=0.05\)，模型与原型为同一流体（水），
  重力起主要作用。（原卷影印见本册附录）}
\end{ti}

\begin{tishi}
\textbf{第 1 题（计算）：}原卷图中除 \(H=3\,\mathrm{m}\)、\(\alpha=45^\circ\)、\(B=1\,\mathrm{m}\) 之外的
几何细节（圆弧半径 \(r\) 与闸门顶缘 \(b\) 的位置、铰链 \(C\) 相对闸底的高程、
水位是否恰好齐闸顶）OCR 未给出，此处\textbf{只按影印可见内容描述、未增补任何数值}，
请对照原卷影印核实后再作答。

\textbf{第 2 题（计算）：}原卷 OCR 作"水从口水池"，
据影印应为"水从\textbf{敞}口水池"，此处已按影印改正。
另原卷图 5 中 \(M\) 点所在管段距出口的\textbf{具体长度（或 \(M\) 点高程）未标出}，
求 \(M\) 点压强时需按题图几何确定 \(M\) 点高程，请对照原卷影印核实。

\textbf{第 3 题（计算）：}原卷 OCR 把题干中的 \(\theta\) 拾作"\(0\)"，
据影印为 \(\theta=30^\circ\)，此处已按影印改正；题干、图中的数值一律照原卷、未改。

\textbf{第 5 题（计算）：}原卷 OCR 把水的运动黏度拾作"\(10^\circ\,\mathrm{m^2/s}\)"、
密度拾作"\(10\,\mathrm{kg/m^3}\)"，据影印分别为
\(\nu=10^{-6}\,\mathrm{m^2/s}\)、\(\rho=10^3\,\mathrm{kg/m^3}\)，此处已按影印改正；
\(\Rey_{x,\mathrm{cr}}=5\times10^5\) 与四个边界层公式照原卷转写。
\underline{数值本身未作任何改动。}

\textbf{第 6 题（计算）：}原卷 OCR 作"长度比尺 \(k=0.05\)"，
据影印为 \(k_l=0.05\)（长度比尺），此处已按影印补下标 \(l\)，数值 \(0.05\) 未改。
题设未说明模型与原型的流体是否相同，
按闸下泄流实验的常规取\textbf{模型与原型为同一流体（水）}。
\end{tishi}
